% ============================================================
% main.tex — クラス替え最適化 授業プリント
% エンジン: LuaLaTeX / ビルド: make (→ out/main.pdf)
% ============================================================
\documentclass[a4paper, 11pt, fleqn]{ltjsarticle}
% ============================================================
% preamble.tex — 授業プリント共通設定
% エンジン: LuaLaTeX + luatexja
% ============================================================

% --- 日本語対応 ---
\usepackage{luatexja}
\usepackage{luatexja-fontspec}

% --- ページレイアウト ---
\usepackage[
  a4paper,
  top=20mm, bottom=20mm,
  left=20mm, right=20mm
]{geometry}

% --- 数式 ---
\usepackage{amsmath}    % 基本数式環境
\usepackage{amssymb}    % 数学記号
\usepackage{amsfonts}   % 数学フォント
\usepackage{physics}    % 物理記法 (\dv, \pdv, \abs, \qty, \vec など)

% --- 単位 ---
\usepackage{siunitx}
\sisetup{
  inter-unit-product = \ensuremath{{}\cdot{}},  % 単位間の区切り
  output-decimal-marker = {.},
  per-mode = symbol
}

% --- グラフ・図 ---
\usepackage{tikz}
\usepackage{pgfplots}
\pgfplotsset{compat=1.18}
\usetikzlibrary{arrows.meta, calc, positioning, decorations.pathmorphing}

% --- その他 ---
\usepackage{graphicx}      % 画像挿入
\usepackage{float}         % 図の位置指定 [H]
\usepackage{booktabs}      % きれいな表 (\toprule, \midrule, \bottomrule)
\usepackage{multicol}      % 段組み
\usepackage{enumitem}      % 箇条書きカスタマイズ
\usepackage{tcolorbox}     % カラーボックス (解答欄など)
\tcbuselibrary{most}
\usepackage{mdframed}      % 枠囲み
\usepackage{xcolor}        % 色

% --- ハイパーリンク (最後に読み込む) ---
\usepackage[hidelinks]{hyperref}

% ============================================================
% カスタムコマンド
% ============================================================

% 解答欄ボックス
\newtcolorbox{answerbox}[1][解答]{
  colback=gray!5,
  colframe=gray!50,
  fonttitle=\bfseries,
  title=#1,
  sharp corners,
  boxrule=0.5pt
}

% 問題番号付き問題環境
\newcounter{mondai}
\newcommand{\mondai}[1][]{%
  \stepcounter{mondai}%
  \medskip
  \noindent\textbf{問題 \themondai}\if\relax\detokenize{#1}\relax\else\ (#1)\fi\quad
}

% 物理でよく使うショートハンド
\newcommand{\kgms}{\,\unit{kg.m/s^2}}
\newcommand{\ms}{\,\unit{m/s}}
\newcommand{\mss}{\,\unit{m/s^2}}

% ヘッダー・フッター設定
\usepackage{fancyhdr}
\pagestyle{fancy}
\fancyhf{}
\lhead{\small \textsf{\@title}}
\rhead{\small \textsf{\@date}}
\cfoot{\small \thepage}
\renewcommand{\headrulewidth}{0.4pt}


\title{クラス替えを数学で解く\\
       \large — 整数計画法入門 —}
\date{\today}

\begin{document}
\maketitle
\thispagestyle{fancy}

\begin{tcolorbox}[colback=yellow!10, colframe=orange!60,
                  title=\textbf{⚠️ このドキュメントについて}]
  本資料は \textbf{Claude Code（Anthropic）} を用いた開発環境のテストとして
  自動生成されたサンプルです。教育目的のデモコンテンツであり、
  実際の授業資料ではありません。\\[4pt]
  \texttt{https://github.com/phys-ken/class-assignment-opt}
\end{tcolorbox}

\bigskip

\begin{tcolorbox}[colback=blue!5, colframe=blue!40, title=\textbf{学習目標}]
  \begin{itemize}
    \item 「最適化」とは何かを理解する
    \item 0-1 変数を使って問題を数式で表現できる
    \item 制約条件と目的関数の意味を説明できる
  \end{itemize}
\end{tcolorbox}

\bigskip

% ============================================================
\section{問題設定}
% ============================================================

全校生徒 $n = 200$ 人を $k = 5$ クラスに割り当てる。
各クラスは 38〜42 人、男女比はそれぞれ 18〜22 人とする。

クラス替えの組み合わせ総数はおよそ

\[
  \binom{200}{40}\binom{160}{40}\binom{120}{40}\binom{80}{40} \approx 10^{130}
\]

であり、全探索は不可能である。そこで\textbf{整数計画法}を用いる。

% ============================================================
\section{数理最適化の3要素}
% ============================================================

\begin{tcolorbox}[colback=gray!5, colframe=gray!40]
  数理最適化 = \textbf{変数} + \textbf{制約条件} + \textbf{目的関数}
\end{tcolorbox}

\subsection{決定変数}

\[
  x_{i,j} \in \{0, 1\}, \quad i = 1, \ldots, 200, \quad j = 1, \ldots, 5
\]

$x_{i,j} = 1$ のとき「生徒 $i$ はクラス $j$ に所属する」を意味する。
変数の総数は $200 \times 5 = 1000$ 個。

\subsection{制約条件}

\begin{enumerate}[label=\textbf{制約\arabic*.}]

  \item \textbf{各生徒は必ず 1 クラスに所属}
  \[
    \sum_{j=1}^{5} x_{i,j} = 1 \qquad (\forall\, i = 1, \ldots, 200)
  \]

  \item \textbf{クラスの人数バランス}
  \[
    38 \leq \sum_{i=1}^{200} x_{i,j} \leq 42 \qquad (\forall\, j = 1, \ldots, 5)
  \]

  \item \textbf{男女比のバランス}（男子の集合 $M$、女子の集合 $F$）
  \[
    18 \leq \sum_{i \in M} x_{i,j} \leq 22, \quad
    18 \leq \sum_{i \in F} x_{i,j} \leq 22 \qquad (\forall\, j)
  \]

  \item \textbf{仲の悪いペア $(a, b)$ を別クラスに（任意）}
  \[
    x_{a,j} + x_{b,j} \leq 1 \qquad (\forall\, j = 1, \ldots, 5)
  \]

\end{enumerate}

\subsection{目的関数}

今回は「制約をすべて満たす解を 1 つ見つける」ことを目標とする
（実行可能解の探索）。目的関数は設定しない。

より発展的な問題では、たとえば「避けたいペアが同じクラスになる回数を最小化」
する目的関数を追加できる。

% ============================================================
\section{OR-tools による求解}
% ============================================================

Google が開発した OR-tools の CP-SAT ソルバーを用いる。

\begin{verbatim}
from ortools.sat.python import cp_model

model = cp_model.CpModel()

# 決定変数
x = [[model.NewBoolVar(f'x_{i}_{j}')
      for j in range(5)] for i in range(200)]

# 制約1: 各生徒は1クラスに所属
for i in range(200):
    model.AddExactlyOne(x[i][j] for j in range(5))

# 制約2: 人数バランス
for j in range(5):
    model.Add(sum(x[i][j] for i in range(200)) >= 38)
    model.Add(sum(x[i][j] for i in range(200)) <= 42)

# 求解
solver = cp_model.CpSolver()
status = solver.Solve(model)
\end{verbatim}

計算時間は通常 \textbf{1 秒以内}。

% ============================================================
\section{考察}
% ============================================================

\mondai[基本]
制約条件① $\displaystyle\sum_{j=1}^{5} x_{i,j} = 1$ が必要な理由を説明せよ。

\begin{answerbox}
  \vspace{2.5cm}
\end{answerbox}

\mondai[発展]
「昨年同じクラスだった生徒ペアの集合 $Q$ を、できるだけ別クラスにしたい」
という条件を制約または目的関数として定式化せよ。

\begin{answerbox}
  \vspace{3.5cm}
\end{answerbox}

% ============================================================
\section*{まとめ}
% ============================================================

\begin{center}
\begin{tabular}{ll}
  \toprule
  用語 & 意味 \\
  \midrule
  整数計画法 & 変数が整数（特に 0-1）に限定された最適化問題 \\
  決定変数   & 値を決めるべき未知数 $x_{i,j} \in \{0,1\}$ \\
  制約条件   & 解が満たすべき等式・不等式 \\
  目的関数   & 最小化・最大化したい量 \\
  CP-SAT     & Google OR-tools の制約プログラミングソルバー \\
  \bottomrule
\end{tabular}
\end{center}

\end{document}
